% ECONOMICS_PROBLEM: {"id": "D44-622", "title": "Trust in complex auction rules", "jel_code": "D44", "source_id": "OP-0295"}
% !TeX program = xelatex
\documentclass[11pt,a4paper]{article}
\usepackage[margin=23mm]{geometry}
\usepackage{fontspec}
\setmainfont{Latin Modern Roman}
\usepackage{amsmath,amssymb,mathtools}
\usepackage{xcolor,enumitem,booktabs,longtable,array,xurl,hyperref}
\definecolor{muted}{HTML}{53636C}
\hypersetup{colorlinks=true,urlcolor=blue,linkcolor=blue}
\setlength{\parindent}{0pt}
\setlength{\parskip}{5pt}
\newcommand{\R}{\mathbb R}
\newcommand{\E}{\mathbb E}
\newcommand{\N}{\mathbb N}
\newcommand{\ind}{\mathbf 1}
\newcommand{\Var}{\operatorname{Var}}
\newcommand{\Cov}{\operatorname{Cov}}
\newcommand{\supp}{\operatorname{supp}}
\DeclareMathOperator*{\argmin}{arg\,min}
\DeclareMathOperator*{\argmax}{arg\,max}
\newcommand{\entrymeta}[3]{{\sffamily\small\color{muted}\textbf{\texttt{#1}}\quad|\quad #2\par #3\par}}
\newcommand{\Problemref}[1]{\texttt{#1}}
\newcommand{\JELref}[2]{\texttt{#2} (JEL \texttt{#1})}
\begin{document}
\section*{Trust in complex auction rules}
\textbf{Problem ID:} \texttt{D44-622}\quad \textbf{JEL:} \texttt{D44}\par
% BEGIN ECONOMICS STATEMENT
\label{problem:OP-0295}
{\sffamily\small\color{muted}Original subject group: Mechanism design and market design\par}
\label{definitions:problem:30007577}
\textbf{Audit classification:} Published research agenda. The verified 2024 survey §4.3 discusses simplicity and credibility. The group-randomized causal attribution problem is editorial.
\subsection*{Definitions and precise research target}
\textbf{Model and research target.} Fix a combinatorial auction with goods \(G\), allocation rule \(x\), payment rule \(p\), and bidder values drawn from a stated population distribution. Randomize groups across interfaces \(Z\) that differ in explanations, rule verification, or commitment evidence while implementing identical economic allocation and payment rules. Let \(Y_i(Z)\) denote bidder participation, \(L_i(Z)\) the loss relative to the bidder's best response under the observed auction environment, and \(W(Z)\) total allocative surplus. Elicit perceived probabilities \(b_i(Z)\) that the auctioneer follows the announced rule, using an incentivized method with specified measurement error. Estimate
\[
 E[Y_i(1)-Y_i(0)],\quad E[L_i(1)-L_i(0)],\quad E[W(1)-W(0)].
\]
Determine how the treatment works through perceived credibility, comprehension, and participation costs. Any mediation decomposition must state its additional identification assumptions; randomization of an explanation alone does not identify the separate causal effect of trust. Include auctioneer-side deviations in a separate prespecified design study, with deviations classified by detectability and the auctioneer's revenue incentives. Evaluate effects on experienced as well as inexperienced bidders and test persistence across auctions.

\emph{Definitions and maintained assumptions (editorial unless a source definition is named).} Fix bidder type law, action spaces, allocation and payment maps, number of rounds and what each bidder observes. Treatments change only explanations or commitment evidence in the maintained economic protocol; if cryptographic auditing changes feasible auctioneer deviations, that is a distinct mechanism treatment. Define loss \(L_i(Z)=\sup_{a_i'}\mathbb E[u_i(a_i',a_{-i})\mid\mathcal I_i,Z]-\mathbb E[u_i(a_i,a_{-i})\mid\mathcal I_i,Z]\), using a stated response belief and the same information for chosen and comparator actions. Ex-post hindsight loss is another estimand. \(W(Z)\) is total value of allocated bundles minus real participation/implementation costs, with transfers counted only under a stated social welfare rule. Beliefs \(b_i\in[0,1]\) are forecasts of an auditable rule-following event. Group assignment identifies total treatment effects despite within-auction interference. Separating credibility from comprehension requires independent manipulations or a justified mediation model; conditioning on elicited beliefs alone does not identify a trust mechanism.
\subsection*{Variants and assumption changes}
\begin{enumerate}
\item \textbf{Explanation factorial (editorial).} Independently vary plain-language explanation and verification evidence while implementing the same committed rules. Estimate total and interaction effects on participation and bidding loss.
\item \textbf{Auctioneer incentives (source-related editorial).} Permit a defined set of deviations and define credibility as following announced rules being a best response against bidder behavior. Compare informational reassurance with devices that actually alter those incentives.
\end{enumerate}
\subsection*{References and source audit}
\begin{enumerate}
\item §4.3 pp.23–25 in author manuscript. Supports the stated research area; exact displayed specification is editorial. Original LSE fetch failed; full manuscript recovered from Harvard author page. \url{https://parkes.seas.harvard.edu/sites/g/files/omnuum8711/files/parkes/files/palacios-huerta.pdf}
\end{enumerate}
\begin{enumerate}\raggedright
\item\hypertarget{definitions:source:30007577:1}{} Ignacio Palacios-Huerta; David C. Parkes; Richard Steinberg. 2024. \emph{Combinatorial Auctions in Practice}. §4.3, printed pp.23–25, Harvard author manuscript; published JEL62(2),517–553 (2024).
\url{https://parkes.seas.harvard.edu/sites/g/files/omnuum8711/files/parkes/files/palacios-huerta.pdf}.
DOI: \url{https://doi.org/10.1257/jel.20221679}.
\end{enumerate}

\subsection*{Common conventions: Source mathematical conventions}

These conventions apply to each entry unless explicitly replaced.

\textbf{Models and identification.} A primitive model \(m\) specifies all probability laws, feasible choices, information, equilibrium and measurement rules used by its target. The admissible class \(\mathcal M\) is fixed independently of outcomes. If \(P_O(m)\) is its observable probability law and \(\tau(m)\) the target, its identified set at observable law \(P\) is
\[
 \mathcal I_\tau(P)=\{\tau(m):m\in\mathcal M,\ P_O(m)=P\}.
\]
Point identification means that this set is a singleton. An empty set signals incompatibility of data and assumptions. Coordinate bounds need not describe the joint set. Bounds are sharp only if every retained target is attainable or arbitrarily approachable by compatible models. Finite bounds require explicit restrictions on unobserved quantities; unrestricted confounding, hidden wealth, extrapolation or arbitrary equilibrium selection can yield uninformative sets.

\textbf{Causal interventions.} A potential outcome \(Y_i(p)\) is the outcome under a complete policy regime \(p\), including timing, financing, information and market spillovers. The target population and units are fixed before treatment. With interference, use the complete assignment vector or a justified exposure mapping; individual no-interference notation alone cannot identify an economy-wide intervention. A difference of mean potential outcomes does not require the cross-world joint distribution. Conditioning on post-treatment migration, survival, enrollment or employment changes the estimand and needs separate justification.

\textbf{Statistical statements.} An estimator, confidence set, forecast or test specifies a sampling law, model class and loss/coverage criterion. Uniform \((1-\alpha)\)-coverage means \(\inf_{m\in\mathcal M}\Pr_m\{\tau(m)\in C_n\}\geq1-\alpha\), or an explicitly asymptotic version. The whole parameter space trivially has coverage, so informativeness requires a stated width, risk or power objective. A forecasting improvement is relative to a named benchmark, information set and evaluation population.

\textbf{Equilibrium and optimization.} An equilibrium comprises feasible plans, optimality, beliefs where needed, budget constraints and all market/resource clearing. The primitives or a stated benchmark define each correspondence \(\mathcal E(p;m)\). Nonemptiness is assumed or proved, never inferred just from compact policy sets. A selected-equilibrium target uses a declared selection rule; a robust target quantifies over all permitted equilibria. Use suprema or infima unless attainment is established. An \(\varepsilon\)-optimal policy is feasible and within \(\varepsilon\) of the finite optimum value. Report an empty feasible set separately.

\textbf{Welfare and accounting.} Normative utility, interpersonal normalization, social weights, numeraire, discounting and the use of public receipts are primitives. Behavioral utility need not coincide with the welfare criterion. Household welfare includes ownership income and rents once; adding profits again to compensating variation generally double-counts them. A monetary equivalent is defined by an explicit transfer and price convention, with monotonicity and range conditions. Average effects, mechanism contributions, transfers and welfare losses are different targets. Infinite sums require convergence and dynamic feasible plans require terminal or no-Ponzi conditions.

\textbf{Source versions.} A working-paper date, revision date and journal publication year are distinguished. A DOI for a journal version is not automatically the DOI of an earlier manuscript. Locators refer to the stated version and printed pages where specified. A metadata-only check does not verify a claim about a full-text conclusion. Every entry's source subsection narrows the provenance claim accordingly.
% END ECONOMICS STATEMENT
\end{document}
